% 05_body.tex — 산출물 5/5 프리모템 기록과 초기 리스크 목록 (빈 템플릿 / 완성 예시 공용 본문)
% 드라이버: 05_프리모템리스크_빈템플릿.tex (\wbblanktrue) / 05_프리모템리스크_완성예시.tex (\wbblankfalse)
% 내용 출처: PM트랙/Day1_워크북.md §5.3(항목 가이드) §5.4(빈 템플릿) §5.5(온담 P1 완성 예시본)
\documentclass[10pt,a4paper]{article}
\usepackage[a4paper,margin=16mm,top=14mm,bottom=20mm]{geometry}
\def\DocNum{5}
\def\DocTitle{프리모템 기록과 초기 리스크 목록}
\def\DocSub{Pre-mortem Record \& Initial Risk List}
\def\DocVariant{\B{빈 템플릿}{온담 P1 완성 예시}}
\usepackage{wbstyle}
\begin{document}

\docheader
 {\Bg{[정식 명칭과 코드]}{차세대 주문·재고 통합 플랫폼 구축 (ONE ONDAM P1)}}
 {\Bg{[판정 시점: 게이트·날짜]}{초판 / 판정 시점 G0+3일 2026-01-08}}
 {\Bg{[진행자 직책·실명]}{P1 PM (발주사 세션 진행)\rd{7mm}}}
 {\Bg{[갱신 책임자 — RACI 26행 A]}{P1 PM (RACI \#26 A: 게이트마다 갱신)\rd{7mm}}}

% ------------------------------------------------------------------ 1
\secbar{1. 세션 정보}
\guide{언제 누가, 어느 시점의 지식으로 했는가. 참가자는 역할로 적어도 된다(개별 사유의 익명성이 전제). 수행사 참여 여부와 세션 분리 여부를 적는다. 임원 참석은 발언을 위축시키고, 운영·현업이 없으면 기술 리스크만 나온다.}
{\footnotesize
\begin{tabularx}{\linewidth}{|L{30mm}|L{26mm}|Y|L{30mm}|L{28mm}|}\hline
\hd{일시} & \hd{진행자} & \hd{참가자 구성(역할·인원)} & \hd{수행사 포함 여부 / 세션 분리} & \hd{지식 기준 시점} \\ \hline
\ifwbblank
\rd{20mm}\g{[날짜, 세션별]} & \g{[직책]} & \g{[발주사·수행사·현업·운영·준법 — 역할과 인원]} & \g{[포함 / 분리 실시]} & \g{[헌장 v, BC v, 정본 날짜]} \\ \hline
\else
2026-01-08 (발주사 세션), 2026-01-09 (합동 세션) & P1 PM (발주사 세션), 프로그램 PMO장 (합동 세션) & 발주사 10명: P1 PM, P1 PMO 3, IT본부 개발 1·운영 1, 영업본부 매장운영 1, 물류본부 센터운영 1, 서비스운영팀 1, 준법지원팀 1 & 합동 세션에 이재현 PM 외 3명 참여 (발주사 세션과 분리) & 헌장 v1.0, 비즈니스 케이스 요약 v0.9, 회사 정본 2025-12 \\ \hline
\fi
\end{tabularx}}

% ------------------------------------------------------------------ 2
\secbar[110mm]{2. 실패 선언문 (Failure Statement — prospective hindsight)}
\guide{실패한 미래의 어느 날을 특정하고 그 모습을 \textbf{완료된 사실(과거형)}로 3~5문장, 날짜·숫자 포함. 일정 실패 / 편익 실패 / 중단 세 종류를 쓰면 사유의 다양성이 커진다. “실패할 수 있다”가 아니라 “실패했다”. 참가자에게 읽어 준 뒤 개별 작성을 시작한다.}
\begin{fieldbox}
\textbf{F1 (일정 실패)} \guide{“[날짜], [무슨 일이 있었다].”}
\B{\lines{3}{8.5mm}}{“2027년 8월이다. P1은 2027-02-26 컷오버 창을 놓쳤다. 다음 창은 성수기와 IPO 예비심사 준비가 겹쳐 2027년 10월로 밀렸고, 예비심사 청구 서류에는 차세대 플랫폼이 ‘구축 중’으로 기재되었다. 대표이사는 착수 회의에서 ‘가능하면 더 앞당기라’고 했었다.”}
\end{fieldbox}
\begin{fieldbox}
\textbf{F2 (편익 실패)} \guide{“[날짜], [\ ].”}
\B{\lines{3}{8.5mm}}{“2028-03-31 G5다. 플랫폼은 2027-02에 가동됐고 관리 성공 기준은 모두 충족했다. 그러나 편익 실현률은 41\%로 판정됐다. 폐기율은 3.9\%에서 3.4\%까지만 내려갔고, 가맹점 판매 데이터는 388점 중 211점만 들어오고 있으며, 앱 주문 비중은 6.2\%에서 8.9\%가 되었다. CFO는 ‘2019년과 무엇이 다른가’라고 물었다.”}
\end{fieldbox}
\begin{fieldbox}
\textbf{F3 (중단)} \guide{“[날짜], [\ ].”}
\B{\lines{3}{8.5mm}}{“2026-11-27 G2다. 프로그램 이사회는 P1 중단 또는 범위 50\% 축소를 의결했다. 1차연도 집행이 148억으로 줄면서 설계 스프린트 인력이 30\% 줄었고, 3PL 협상은 타결되지 않았으며, P2 기계적 준공은 2027-05로 통지됐다. 편익 전망은 손익분기 66\%를 밑돌았다. 종료 보고서의 중단 사유 칸은 이번에는 세 줄보다 길었다.”}
\end{fieldbox}

% ------------------------------------------------------------------ 3
\secbar[80mm]{3. 개별 실패 사유 원문 (Raw Reasons)}
\guide{참가자가 각자 적은 사유를 편집 없이 번호를 붙여 옮긴다. 익명, 과거형(“~할 것 같다”는 다시 쓰게 한다), 한 줄 한 사유. 중복도 그대로 둔다 — 중복 횟수가 우선순위 정보다. 수집 단계에서 토론·삭제·범주화를 하지 않는다.}
\begin{fieldbox}
\ifwbblank
\foreach \i in {1,...,20}{\g{\i.} \rule{0pt}{8mm}\hrulefill\par}
\g{(계속되면 뒷면 또는 별지에)}
\else
\guide{(군집별 대표 원문, 총 61건 중 발췌)}
\begin{enumerate}
\item 가맹점주협의회가 단말 비용과 데이터 용도 문제로 2026-08 부속서 서명을 거부했고, 그때 우리는 이미 가맹 SW를 개발 중이었다.
\item 대성로지스가 2022년과 똑같이 “시스템 부하”를 이유로 준실시간 연동을 거부했고, 물류본부는 그 협상을 P1 일정에 맞춰 시작하지 않았다.
\item P2 설비가 41주로 늦어졌고, 물류 인터페이스 시험을 못 한 채 컷오버를 갈지 말지를 G3 전날까지 아무도 결정하지 못했다.
\item CFO가 1차연도 집행을 148억으로 줄였는데 P1의 어느 항목을 줄일지는 3개월 동안 정해지지 않았고, 라이선스 선약정이 밀려 인프라가 설계 스프린트 끝까지 없었다.
\item 발주사 인력 14명이 겸임이라 요구사항 확정 회의에 절반이 안 왔고, 수행사는 대기 시간을 전부 기록해 변경협의체에 올렸다.
\item OD-POS를 아는 마지막 개발자가 2026년 봄에 퇴사했고, 인터페이스 13본은 아무도 읽을 수 없었다.
\item 요구사항 동결 전날마다 영업본부에서 화면 요구가 40건씩 들어왔고, 범위 허용범위 ±35건을 넘긴 달에 FP 재산정이 걸려 한 달을 잃었다.
\item 자사앱 벤더가 P1 API 연동을 “계약 외”라며 별도 견적 4억을 냈고, 앱 주문 연동이 컷오버에서 빠졌다.
\item 개인정보 처리항목 142개 중 온담푸드시스템(별도 법인) 몫 31개의 처리 주체가 2026-09-11 직전까지 안 정해졌다.
\item 데이터 전환 정본 판정 뒤에 상품 마스터 코드 체계가 바뀌었고, 610개 테이블을 다시 매핑하느라 리허설 3회차가 밀렸다.
\item 컷오버 주말에 가맹점 단말 중 2013년 이전 기종 90여 대가 펌웨어를 못 받았고, 월요일 아침에 그 매장들이 전화·팩스로 돌아갔다.
\item 박준영 팀장이 G4에서 SAC 미해소 항목 23건을 들어 인수를 거부했고, P1은 종료되지 못한 채 5개월간 하자보수 체제로 남았다.
\item 노조가 이천센터 실사 협조를 거부해 재고 데이터 검증 표본을 못 뽑았다.
\item 편익 측정 책임자로 지정된 본부장이 “그 숫자는 시스템 탓이 아니라 경기 탓”이라고 했고, 분모가 다른 두 숫자가 이사회에 동시에 올라갔다.
\item 감리가 G2에서 RTM 커버리지 미달을 시정조치로 지적했는데 대응에 6주가 걸렸다.
\end{enumerate}
\fi
\end{fieldbox}

% ------------------------------------------------------------------ 4
\secbar[120mm]{4. 군집화와 우선순위 (Clustering \& Prioritization)}
\guide{원문을 유사 원인별로 묶고 언급 수와 참가자 투표(1인 3표)로 순위. 군집명은 범주(“일정”)가 아니라 \textbf{원인}(“가맹 부속서 협상 지연”)으로. 군집을 너무 크게 잡지 말고(“외부 요인”), 득표가 적어도 치명적인 사유(원 개발자 1명 이탈)를 탈락시키지 말 것.}
\B{}{\small (10명, 1인 3표)\par}
{\footnotesize
\begin{tabularx}{\linewidth}{|L{12mm}|Y|L{22mm}|L{16mm}|L{14mm}|L{14mm}|}\hline
\hd{군집} & \hd{군집명 (원인으로)} & \hd{대표 원문 \#} & \hd{언급 수} & \hd{득표} & \hd{순위} \\ \hline
\ifwbblank
\rd{8mm}C1 & \g{[원인으로 쓴 군집명]} & \g{[원문 번호]} & & & \\ \hline
\rd{8mm}C2 & & & & & \\ \hline
\rd{8mm}C3 & & & & & \\ \hline
\rd{8mm}C4 & & & & & \\ \hline
\rd{8mm}C5 & & & & & \\ \hline
\rd{8mm}C6 & & & & & \\ \hline
\rd{8mm}C7 & & & & & \\ \hline
\rd{8mm}C8 & & & & & \\ \hline
\rd{8mm}C9 & & & & & \\ \hline
\rd{8mm}C10 & & & & & \\ \hline
\rd{8mm}C11 & & & & & \\ \hline
\rd{8mm}C12 & & & & & \\ \hline
\else
C1 & P2 리드타임 41주 통지 — 물류 인터페이스·컷오버 결합 & 3 & 9 & 8 & 1 \\ \hline
C2 & 가맹 부속서 협상(2026-08-14) 지연·결렬 & 1, 11 & 8 & 7 & 2 \\ \hline
C3 & 3PL 데이터 개방·연동 거부 재발 & 2 & 7 & 5 & 3 \\ \hline
C4 & 1차연도 집행 148억 축소와 P1 배분 미결정 & 4 & 6 & 5 & 3 \\ \hline
C5 & 발주사 겸임 인력 투입률·미결정 대기 → 변경대가 & 5 & 6 & 4 & 5 \\ \hline
C6 & OD-POS 지식 보유자 1명 의존 & 6 & 3 & 4 & 5 \\ \hline
C7 & 요구사항 동결 직전 집중 → 범위 허용범위 초과 & 7 & 5 & 2 & 7 \\ \hline
C8 & 자사앱 벤더 소스 보유·연동 계약 외 주장 & 8 & 4 & 3 & 7 \\ \hline
C9 & 별도 법인(온담푸드시스템) 규제·계약 주체 미확정 & 9 & 3 & 2 & 9 \\ \hline
C10 & 데이터 전환 정본 판정 후 마스터 변경 & 10 & 4 & 3 & 7 \\ \hline
C11 & SAC 미합의로 G4 인수 거부 & 12 & 3 & 2 & 9 \\ \hline
C12 & 편익 측정 분모·책임자 분쟁 / 노조 실사 비협조 & 13, 14 & 4 & 3 & 7 \\ \hline
\fi
\end{tabularx}}

% ------------------------------------------------------------------ 5 (가로)
\landscapeon
\secbar{5. 초기 리스크 목록 (Initial Risk List)}
\guide{군집을 리스크 서술로 전환. 서술은 반드시 3부 — \textbf{원인}(현재 사실) → \textbf{사건}(불확실한 미래) → \textbf{영향}(목표에의 결과). 이미 발생한 사실(OD-POS OS 지원 종료)은 리스크가 아니다. 확률·영향은 상중하까지(숫자는 Day2). 소유자는 가장 잘 관리할 수 있는 사람 — PM이 아닌 경우가 많다. 대응은 방향만(회피·전가·완화·수용 / 기회는 활용·공유·증대·수용). 기회(O) 1~2건 포함. 헌장 항목 10의 상위 리스크가 이 목록의 상위와 일치해야 한다.}
{\scriptsize
\begin{longtable}{|L{9mm}|L{62mm}|L{15mm}|L{15mm}|L{8mm}|L{8mm}|L{20mm}|L{20mm}|L{40mm}|L{24mm}|L{12mm}|}\hline
\hd{ID} & \hd{원인 → 사건 → 영향} & \hd{범주} & \hd{출처} & \hd{확률} & \hd{영향} & \hd{근접성} & \hd{소유자} & \hd{초기 대응 방향} & \hd{트리거 신호} & \hd{다음 검토} \\ \hline \endhead
\ifwbblank
\rd{15mm}R-01 & \g{[현재 사실] → [불확실한 미래] → [목표에의 결과]} & \g{[외부의존/계약/기술/조직/규제/재무]} & \g{[군집\#/가정\#]} & \g{상중하} & \g{상중하} & \g{[언제 발현 가능]} & \g{[실명]} & \g{[회피/전가/완화/수용 — 방향만]} & \g{[무엇을 보면 발현 신호인가]} & \g{[게이트]} \\ \hline
\rd{15mm}R-02 & & & & & & & & & & \\ \hline
\rd{15mm}R-03 & & & & & & & & & & \\ \hline
\rd{15mm}R-04 & & & & & & & & & & \\ \hline
\rd{15mm}R-05 & & & & & & & & & & \\ \hline
\rd{15mm}R-06 & & & & & & & & & & \\ \hline
\rd{15mm}R-07 & & & & & & & & & & \\ \hline
\rd{15mm}R-08 & & & & & & & & & & \\ \hline
\rd{15mm}R-09 & & & & & & & & & & \\ \hline
\rd{15mm}R-10 & & & & & & & & & & \\ \hline
\rd{15mm}\textbf{O-01} & \g{(기회) [\ ] → [\ ] → [\ ]} & & & & & & & \g{[활용/공유/증대/수용]} & & \\ \hline
\rd{15mm}\textbf{O-02} & \g{(기회)} & & & & & & & & & \\ \hline
\else
R-01 & P2 설비 벤더가 리드타임 34→41주 가능성을 통지한 상태다 → MC가 2027-02-12를 넘겨 2027-04 이후가 된다 → 물류 인터페이스 통합시험을 G3 전에 못 하고, 컷오버를 매장·주문 부분만 먼저 하거나 창을 이탈한다 & 외부 의존 & C1 / 헌장 가정 5 & 상 & 상 & T0 2026-04-24에 확정 & 프로그램 PMO장 (P1측: P1 PM) & 완화 — 컷오버를 “매장·주문 1차 / 물류 2차”로 분리 가능한 설계를 G1까지 옵션으로 준비, 프로그램 이사회 결정 요청 & P2 T0 시 리드타임 확정 통보, 2026-10 P2 진도 & G1 \\ \hline
R-02 & 가맹 388점의 단말 비용·데이터 용도 합의가 없고 협의회가 이미 문제를 제기했다 → 2026-08-14 부속서 개정이 지연되거나 일부 매장이 거부한다 → 가맹 SW 배포 범위 축소, 가맹 판매 데이터 미수집, 편익 B-02·B-04 가맹 부분 미실현 & 조직·외부 & C2 / 가정 2 & 상 & 상 & 2026-06~08 & 오정근 & 완화 — 가맹 포털(이익)과 POS SW(부담)를 분리 릴리스, 데이터 용도 명세서와 단말 호환 조사(2013년 이전 기종 수) 2026-04까지 & 협의회 서면 질의, 동의율 2026-07 중간 집계 & 2026-06, G1 \\ \hline
R-03 & 대성로지스 WMS는 벤더 소유이며 2022년 데이터 개방 거부 이력이 있다 → 준실시간 연동 계약 변경이 2026-06까지 타결되지 않는다 → 센터 재고가 야간 배치로 남아 B-02 목표가 매장 재고에 한정 & 외부 의존·계약 & C3 / 가정 1 & 중 & 상 & 2026-03~06 & 한동석 & 완화+수용 — 야간 배치 2회 폴백 설계, 연동 부하 산정서 제공으로 “시스템 부하” 논거 선제 대응 & 협상 착수 여부 2026-03, 대성 서면 회신 & G1 \\ \hline
R-04 & CFO가 1차연도 집행 148억 축소를 요구하고 있다 → 확정 시 P1 2026년 배분이 줄고 어느 항목을 줄일지 지연된다 → 라이선스 선약정·인프라 지연, 설계 스프린트 인력 축소, 컷오버 창 이탈 & 재무 & C4 / 헌장 리스크 4 & 상 & 상 & 2026-02 자금계획 확정 & 정미란 (대응안: P1 PM) & 완화 — 삭감 시나리오별 P1 항목 축소안(이월 vs 총액 삭감) 2026-01 말까지 선제 제출 & 재무본부 자금계획 2026 확정 & 2026-02 \\ \hline
R-05 & 발주사 투입 14명이 대부분 겸임이고 수행사는 대기 시간을 변경대가로 기록하겠다고 밝혔다 → 요구사항 확정·검토 회의 참석률 저하 → 결정 지연 대기 비용 청구, 요구사항 베이스라인 v1.0 지연 & 조직·계약 & C5 / 가정 3 & 상 & 중 & 2026-01~03 설계 스프린트 & 박세연 & 완화 — 결정 요청 응답 SLA 5영업일 명문화(RACI), 겸임 인력 투입률 월 측정, 대체 결정권자 지정 & 회의 참석률 70\% 미만 2주 연속 & G1 \\ \hline
R-06 & OD-POS 원 개발자 3명 중 1명만 재직 중이고 인터페이스 13본은 문서가 없다 → 그 1명이 이탈하거나 병행 운영 장애 대응에 묶인다 → 13본 분석·데이터 전환 매핑 불가, G2 정본 판정 지연 & 기술·조직 & C6 & 중 & 상 & 상시, 특히 2026-Q1 & 박세연 & 완화 — 13본 문서화를 설계 스프린트 1~3에 전진 배치, 지식 이전 세션 기록, 인사본부와 유지 방안 협의 & 퇴직 의사 표명, 문서화 진도 2026-03 50\% 미만 & G1 \\ \hline
R-07 & 요구사항 동결이 월 1회이고 영업본부가 SRS 570건을 늘리려 한다 → 동결 직전 주에 요구가 집중해 월 변동이 ±35건을 넘는다 → FP 재산정 트리거, 기준선 재확인, 설계 지연 & 범위·계약 & C7 & 상 & 중 & 매월 마지막 수요일 & P1 PM & 완화 — 동결 2주 전 요구 접수 마감, 우선순위(MoSCoW, Must ≤ 60\%) 규칙 헌장 우선순위와 연동 & 월 접수 건수 추이 & 매월 \\ \hline
R-08 & 자사앱 소스 형상을 외주 벤더가 보유하고 API 연동이 계약 범위에 명시되어 있지 않다 → 벤더가 연동을 계약 외로 주장하고 별도 견적을 낸다 → 앱 주문 연동 지연, B-05·B-01 편익 M+24 이후로 & 계약·외부 & C8 & 중 & 상 & 2026-04 앱 연동 설계 & 박세연 (계약: 오정근) & 전가+완화 — 앱 유지보수 계약에 소스 인수·API 연동 조항 추가(2026-03), 대안으로 P1 예산 외 계약 변경 예비 산정 & 벤더 API 규격 회신 지연 & G1 \\ \hline
R-09 & 온담푸드시스템은 별도 법인이며 개인정보 처리항목 142개 중 공장 몫의 처리 주체가 미확정이다 → 2026-09-11 전까지 위탁 구조가 확정되지 않는다 → 공장 연동을 규제 대응 범위에서 빼거나 P1이 흡수 & 규제 & C9 & 중 & 상 & 2026-06 (규제 버퍼 22영업일 고려) & 최영주 (계약 구조: P1 PM) & 완화 — 별도 법인 위탁 계약 초안 2026-03, P4 규제 아키텍처 확정 시점과 연동 & P4 위탁 구조 확정 여부 2026-05 & G1 \\ \hline
R-10 & 데이터 전환 정본 판정(G2)은 되돌리면 610 테이블 재매핑 3.8억·6주가 든다 → 판정 후 상품 마스터 코드 체계가 바뀐다(신규 브랜드·B2B 요구) → 리허설 3회차 지연, 컷오버 창 이탈 & 기술 & C10 & 중 & 상 & 2026-11-27 이후 & 박세연 & 회피 — G2 전에 마스터 6종 변경 동결 규칙을 CCB 의결, 판정 후 변경은 컷오버 이후 릴리스로 & G2 전 마스터 변경 요청 건수 & G2 \\ \hline
R-11 & SAC가 G1까지 합의되지 않은 상태다 → G4에서 인수자가 미해소 항목을 들어 인수를 거부한다 → P1 종료 불가, 하자보수 체제 장기화 & 조직 & C11 & 중 & 중 & G4 2027-05 & 박준영 (초안: P1 PM) & 완화 — SAC 초안 G1까지 공동 작성, G3에서 SAC 사전 점검 & G1 SAC 미합의 & G1 \\ \hline
R-12 & 편익 분모(순환실사 SKU 69\%)가 불완전하고 책임자 지정이 문서화되지 않았으며 노조가 물류 편익을 인력 감축으로 읽고 있다 → G5에서 분모 분쟁·측정 거부·실사 비협조 → 편익 판정 불가 또는 “미달” & 조직·재무 & C12 / BC 가정 A3 & 중 & 상 & M+3 분모 확정, 실사 2026-Q3 & 오정근·한동석 (분모: 나영선 협의) & 완화 — 편익 책임자 서명 G1, 분모 확장 방안 M+3, P1 편익에 인건비 없음 문서화 후 노조 채널(인사본부) 공유 & 분모 협의 M+3 미완, 실사 협조 거부 & G1 \\ \hline
R-13 & 컷오버는 615개 매장 펌웨어 일괄배포 34시간이며 롤백은 2.1억·6영업일이다 → 가맹 구형 단말이 펌웨어를 받지 못한다 → 일부 매장 수기 복귀, 월요일 장애 콜 폭증 & 기술 & C2(11번 원문) & 중 & 중 & 2027-02-26 & P1 PM (단말 조사: 오정근) & 완화 — 가맹 단말 기종 전수 조사 2026-04, 비호환 기종 대응안(교체 비용 부담 결정은 R-02와 연동), 파일럿 12개점에 구형 기종 포함 & 단말 조사 결과 비호환 10\% 초과 & G1 \\ \hline
\textbf{O-01} & (기회) 폴백으로 남긴 do-minimum(인터페이스 13본 문서화)은 어차피 P1 범위다 → 설계 스프린트 초반에 완료하면 데이터 전환 매핑이 빨라진다 → G2 정본 판정 조기 확보, 리허설 여유 & 기술 & C6 역발상 & 중 & 중 & 2026-Q1 & 박세연 & 활용 — R-06 완화 조치와 통합 & 문서화 진도 & G1 \\ \hline
\textbf{O-02} & (기회) 가맹 포털은 협의회가 환영하는 유일한 부분이다 → 포털을 먼저 릴리스(2026-Q4)하면 부속서 협상의 지렛대가 된다 → R-02 확률 감소, 리드타임 편익 M+12 조기 실현 & 조직 & C2 역발상 & 중 & 상 & 2026-06 협상 전 시연 & 오정근 (P1 PM 지원) & 활용 — 헌장 우선순위 “범위 조정”에서 포털을 이연 불가 항목으로 & 프로토타입 2026-05 & G1 \\ \hline
\fi
\end{longtable}}
\landscapeoff

% ------------------------------------------------------------------ 6
\secbar[90mm]{6. 즉시 조치와 헌장·계획 반영 사항 (Immediate Actions \& Feed-forward)}
\guide{리스크 대응 계획을 기다리지 않고 지금 바꿔야 하는 것 — 헌장의 가정·제약·제외 범위 수정, 이해관계자 등록부 추가, RACI 행 추가. 줄마다 “어느 문서의 어느 항목을 어떻게”와 담당·기한. 프리모템 결과가 리스크 목록에만 남고 헌장은 그대로면, 헌장의 가정은 프리모템이 이미 의심한 것을 계속 ‘참’이라고 말한다.}
\begin{fieldbox}
\ifwbblank
\foreach \i in {1,...,7}{\g{–} \rule{0pt}{9mm}\hrulefill\ \g{(담당: \hspace{16mm}, 기한: \hspace{16mm})}\par}
\else
\begin{itemize}
\item 헌장 항목 5 제외 범위에 “가맹점 단말 HW 교체” 명시 — 반영 완료(v1.0). 담당 P1 PM.
\item 헌장 항목 9 가정에 “자사앱 벤더 연동 2026-12” 추가, 항목 10 리스크에 R-08 반영 — 담당 P1 PM, 2026-01-15 (v1.1).
\item 이해관계자 등록부 S-19 자사앱 벤더를 “위험” 유형으로 등록, 관계 소유자 박세연 — 반영 완료.
\item RACI 12행(앱 소스 인수 조항), 13행(별도 법인 계약 구조) 추가 — 반영 완료.
\item 비즈니스 케이스 요약 가정 A3(분모)에 R-12 연결, 편익 책임자 서명을 G1 통과 조건에 추가 — 담당 P1 PM, 스폰서 협의 2026-01-20.
\item CFO 삭감 시나리오별 P1 축소안(R-04) 2026-01-30까지 스폰서 제출 — 담당 P1 PM·재경팀장.
\item 가맹 단말 기종 전수 조사(R-13) 착수 요청 — 담당 오정근(영업본부 매장운영), 2026-04 완료.
\end{itemize}
\fi
\end{fieldbox}

\end{document}
